%% 第 11 章：广义特征空间及它的主分解
%% 本文件由 tools/md2tex.py 自动生成，请勿直接编辑。

\chapter{广义特征空间及它的主分解}


\begin{jdefn}{广义特征向量}{r:12:1}
若非零向量 $v$ 满足：

\[
(T-\lambda I)^k v=0
\]

对某个正整数 $k$ 成立，则称 $v$ 是属于 $\lambda$ 的广义特征向量。

普通特征向量是 $k=1$ 时的特殊情形。

类比之前的不同的特征值对应的特征向量线性无关，对于广义特征向量，不同的特征值对应的广义特征向量互相线性无关，证明从略。
\end{jdefn}

\begin{jdefn}{广义特征空间}{r:12:2}
定义：

\[
G_\lambda
=
\ker(T-\lambda I)^{\dim V}.
\]

它称为 $\lambda$ 对应的广义特征空间。

实际上，零空间链可能在更早的某一步稳定；使用 $\dim V$ 只是统一而安全的写法。因为零空间链会递增，并且最远在 $\dim V$ 处停止。
\end{jdefn}

\section{广义特征空间与普通特征空间}

\[
E_\lambda
=
\ker(T-\lambda I)
\subseteq
G_\lambda.
\]

普通特征空间只记录被 $T-\lambda I$ 一次消掉的方向。

广义特征空间则记录经过有限次作用后最终被消掉的所有方向。

\begin{jthm}{不同特征值的广义特征空间直和}{r:12:3}
设：

\[
\lambda_1,\ldots,\lambda_s
\]

是 $T$ 的全部不同特征值。若 $V$ 是有限维复向量空间，则：

\[
V
=
G_{\lambda_1}\oplus\cdots\oplus G_{\lambda_s}.
\]

这称为广义特征空间的主分解。
\end{jthm}

\begin{jproof}{主分解的证明思路}
对第一个特征值 $\lambda_1$，令：

\[
N_1=T-\lambda_1I.
\]

对足够大的 $n$，Fitting 分解给出：

\[
V
=
\ker N_1^n
\oplus
\operatorname{range}N_1^n.
\]

其中：

\[
\ker N_1^n=G_{\lambda_1}.
\]

而：

\[
\operatorname{range}N_1^n
\]

对 $T$ 不变，因此可以将 $T$ 限制到这个剩余空间上继续研究。

在该剩余空间中，$\lambda_1$ 已经不再出现为特征值。对下一个特征值重复这一过程，最终得到全部广义特征空间的直和。

Cayley-Hamilton 定理保证特征多项式中的全部因子作用后会消掉整个空间，因此这个过程最终会覆盖整个 $V$。
\end{jproof}

\section{广义特征空间的维数}

在复数域上：

\[
\dim G_\lambda
=
\lambda\text{ 的代数重数}.
\]

而：

\[
\dim E_\lambda
=
\lambda\text{ 的几何重数}.
\]

代数重数表示广义特征空间的总维数；几何重数表示其中普通特征向量所张成空间的维数。
